% !TeX document-id = {ecbf6175-be07-4601-b8eb-d3263d2bff07}
% !TeX TXS-program:compile = txs:///pdflatex/[--shell-escape]
% <<<<<<< HEAD
% \documentclass[accentcolor=3c,landscape,ngerman,presentation,t,usenames,dvipsnames,svgnames,table,colorframetitle=true]{tudabeamer}
% =======

%include guard:
\ifdefined\GlobalAlreadyIncluded
    \expandafter\endinput
\fi

\gdef\GlobalAlreadyIncluded{}
%

\documentclass[accentcolor=3c,landscape,ngerman,presentation,t,usenames,dvipsnames,svgnames,table, aspectratio=169,colorframetitle=true]{tudabeamer}
%>>>>>>> online_rebase

% Template-Modifikationen
\addtobeamertemplate{frametitle}{}{\vspace{-1em}} % mehr Platz vor dem Inhalt

% andere global gemeinsame definitionen
%Includes
\usepackage[ngerman]{babel} %Deutsche Silbentrennung
\usepackage[utf8]{inputenc} %Deutsche Umlaute
\usepackage{float}
\usepackage{graphicx}
\usepackage{minted}
\RequirePackage{csquotes}
\RequirePackage{fontawesome5}

\DeclareGraphicsExtensions{.pdf,.png,.jpg}

\makeatletter
\author{Vorkursteam der Fachschaft Informatik}
\let\Author\@author

% dark mode
\ExplSyntaxOn
\IfDarkModeT{
    \cs_if_exist:NT \setbeamercolor {
        \setbeamercolor*{smallrule}{bg=.}
        \setbeamercolor*{normal~text}{bg=\thepagecolor,fg=.}
        \setbeamercolor*{background~canvas}{parent=normal~text}
        \setbeamercolor*{section~in~toc}{parent=normal~text}
        \setbeamercolor*{subsection~in~toc}{parent=normal~text,fg=.}
        \setbeamercolor*{footline}{parent=normal~text}
        \setbeamercolor{block~title~alerted}{fg=white,bg=white!20!\thepagecolor}
        \setbeamercolor*{block~body}{bg=black!70!gray!98!blue}
        \setbeamercolor*{block~body~alerted}{bg=\thepagecolor}
    }
    \cs_if_exist:NT \setbeamertemplate {
        \setbeamertemplate{subsection~in~toc~shaded}[default][50]
    }
}
\ExplSyntaxOff

% macros
\renewcommand{\arraystretch}{1.2} % Höhe einer Tabellenspalte minimal erhöhen
\newcommand{\N}{{\mathbb N}}
\renewcommand{\code}{\inputminted[]{python}}

\IfDarkModeTF{
    \newmintedfile[pythonfile]{python}{
        fontsize=\small,
        style=native,
        linenos=true,
        numberblanklines=true,
        tabsize=4,
        obeytabs=false,
        breaklines=true,
        autogobble=true,
        encoding="utf8",
        showspaces=false,
        xleftmargin=20pt,
        frame=single,
        framesep=5pt,
    }
    \newmintinline{python}{
        style=native,
        encoding="utf8"
    }
    \newmintinline{kotlin}{
        style=native,
        encoding="utf8"
    }


    \definecolor{codegray}{HTML}{eaf1ff}
    \newminted[bashcode]{awk}{
        escapeinside=||,
        fontsize=\small,
        style=native,
        linenos=true,
        numberblanklines=true,
        tabsize=4,
        obeytabs=false,
        breaklines=true,
        autogobble=true,
        encoding="utf8",
        showspaces=false,
        xleftmargin=20pt,
        frame=single,
        framesep=5pt
    }
}{
    \newmintedfile[pythonfile]{python}{
        fontsize=\small,
        style=friendly,
        linenos=true,
        numberblanklines=true,
        tabsize=4,
        obeytabs=false,
        breaklines=true,
        autogobble=true,
        encoding="utf8",
        showspaces=false,
        xleftmargin=20pt,
        frame=single,
        framesep=5pt,
    }
    \newmintinline{python}{
        style=friendly,
        encoding="utf8"
    }
    \newmintinline{kotlin}{
        style=friendly,
        encoding="utf8"
    }

    \definecolor{codegray}{HTML}{eaf1ff}
    \newminted[bashcode]{awk}{
        escapeinside=||,
        fontsize=\small,
        style=friendly,
        linenos=true,
        numberblanklines=true,
        tabsize=4,
        obeytabs=false,
        breaklines=true,
        autogobble=true,
        encoding="utf8",
        showspaces=false,
        xleftmargin=20pt,
        frame=single,
        framesep=5pt
    }
}

\let\origpythonfile\pythonfile
\renewcommand{\pythonfile}[1]{\pythonfileh{#1}{}}
\newcommand{\pythonfileh}[2]{\origpythonfile[#2]{#1}}

\DeclareDocumentCommand{\kotlinfile}{O{} O{} m}{\inputCode[#1]{minted language=kotlin,#2}{#3}}

\newcommand*{\ditto}{\texttt{\char`\"}}

\newcommand{\shellprefix}{\textcolor{TUDa-3a}{\ttfamily\bfseries \$~}}
\DeclareTCBListing{commandshell}{ O{} O{} }{
    colback=\IfDarkModeTF{black}{black!80},
    colupper=white,
    colframe=TUDa-3a,
    listing only,
    % listing options={style=tcblatex,language=sh},
    listing engine=minted,
    minted style=dracula,
    minted options={
        % linenos=true,
        numbersep=3mm,
        texcl=true,
        autogobble,
        escapeinside=@@,
        breaklines,
        highlightcolor=yellow!50!black,
        #1
    },
    #2,
    % before upper={\textcolor{red}{\small\ttfamily\bfseries root \$> }},
    % every listing line={\textcolor{red}{\small\ttfamily\bfseries root \$> }}
}


%Includes
\usepackage{epstopdf}
\usepackage{wrapfig}
\usepackage{tipa}
\usepackage{tikz}
\usetikzlibrary{calc,shapes,arrows}
%tip: use http://l04.scarfboy.com/coding/phonetic-translation?from=ipa&fromtext=%CB%88pa%C9%AA%CE%B8n%CC%A9&to=tipa
%for converting ipa

\graphicspath{ {./media/} }

\def\shortyear#1{\expandafter\shortyearhelper#1}
\def\shortyearhelper#1#2#3#4{#3#4}

\newcount\NextYear
\NextYear = \year
\advance\NextYear by 1

\newcommand\NextYearShort{\shortyear{\the\NextYear}}

% notes
\usepackage{pgfpages}
\setbeamertemplate{note page}[plain]
%\setbeameroption{show notes on second screen}

% macro for change speaker sign
\newcommand{\changespeaker}{
    \begin{tikzpicture}[line width=.6mm, shorten >= 3pt, shorten <= 3pt]

        \coordinate (c1);
        \coordinate[right of=c1] (c2);

        \draw[rectangle, draw=red!80, fill=red!80, align=center, rounded corners] ($(c1.north west)+(0,-0.3)$) rectangle ($(c2.south east)+(0, 0.3)$) {};
        \draw[->,white] (c1)[bend left] to node[auto] {} (c2);
        \draw[->,white] (c2)[bend left] to node[auto] {} (c1);
    \end{tikzpicture}
}

%Listing-Style pyhon
\title[Programmiervorkurs]{Programmiervorkurs Wintersemester \the\year/\NextYearShort}
\subtitle{{\small der Fachschaft Informatik}}
\logo*{\includegraphics{../globalMedia/bildmarke_ohne_rand}}
\institute{Fachschaft Informatik}
\date{Wintersemester \the\year/\NextYearShort}


% macros
\newcommand{\livecoding}{
    %\ifdefined\StreamSlides
    \begin{frame}
        \frametitle{\insertsectionhead \\  {\small \insertsubsectionhead}}\centering \huge 	\vskip 2cm\textbf{\textcolor{red}{Live-Coding}}
    \end{frame}
    %\fi
}

%\newcommand{\slidehead}{\frametitle{\insertsectionhead \\ {\small \insertsubsectionhead}}\vspace{3mm}}
\newcommand{\slidehead}{\frametitle{\insertsectionhead} \framesubtitle{\insertsubsectionhead}\vspace{3mm}}
\newcommand{\tocslide}{\begin{frame}[t]\frametitle{Inhaltsverzeichnis}\vspace{3mm}{\small\tableofcontents[subsectionstyle=shaded]}\end{frame}}

\newcommand{\nextvid}[2]{
    \ifdefined\StreamSlides
    \else
        \section{Nächstes Video}
        \begin{frame}[t]
            \slidehead
            \begin{block}{Nächstes Video}
                \vspace{0.5cm}
                #1
                \vspace{0.5cm}
            \end{block}
            \ifx\hfuzz#2\hfuzz
                \vspace{2.5cm}
                \else
                {\begin{block}{Bonus Video}
                        \vspace{0.5cm}
                        #2
                        \vspace{0.5cm}
                    \end{block}}
            \fi
            Danke fürs Zuschauen!\\
            Links zu den Folien und Quellen sind in der Videobeschreibung.
        \end{frame}
    \fi
}


\usepackage{verbatim}
\usetikzlibrary{decorations.pathreplacing}
\usetikzlibrary{shapes.misc}


% colors
\definecolor{lightpetrol}{RGB}{\IfDarkModeTF{0,136,119}{0,223,194}}


\begin{document}
\subtitle{Kapitel 2: Spaß mit Daten}
\titlegraphic{
    \begin{columns}
        \begin{column}{6cm}
            \vspace{4.32mm}
            \begin{center}
                {\huge Datentypen}
            \end{center}
            \vspace{0.1cm}
            \begin{figure}
                \centering
                \includeinvertablegraphics[scale=1]{media/su_doku.png}
                \\	\sffamily \tiny Bild: \href{https://xkcd.com/74/}{https://xkcd.com/74/}
            \end{figure}
        \end{column}
    \end{columns}}
\maketitle

\section{Was sind Typen?}
\subsection{Überblick}
\begin{frame}
    \slidehead
    \begin{itemize}
        \item Daten können unterschiedliche \textit{Typen} haben
        \item Wichtige primitive Datentypen:
    \end{itemize}
    \begin{center}
        \begin{tabular}[h]{|c|c|c|}
            \hline
            \textbf{Typ} & \textbf{Wertebereich} & \textbf{Beispiele}                                                                                    \\
            \hline
            Boolean      & Wahrheitswert         & \kotlininline{true}, \kotlininline{false}                                                             \\
            \hline
            Int          & Ganzzahl              & 42, -2                                                                                                \\
            \hline
            Double       & Fließkommazahl        & 42.5, 0.1                                                                                             \\
            \hline
            String       & Text                  & \kotlininline{"test"}, \kotlininline{"Zeilen\numbruch"}                                               \\
            \hline
            Char         & einzelnes Textzeichen & \kotlininline{'a'}, \kotlininline{'€'}, \kotlininline[escapeinside=||]{|\PYG{l+s}{'\PYGZbs{}u263A'}|} \\
            \hline
        \end{tabular}\\
    \end{center}
    {\footnotesize Diese Liste ist unvollständig, es gibt noch weitere Typen.} % Zum Beispiel gibt es einen für Komplexe Zahlen.
    \pause
    \begin{block}{Wichtig}
        Nutze einen Punkt als Dezimaltrenner für Fließkommazahlen!
    \end{block}
\end{frame}

\subsection{Technische Details}
\begin{frame} % Technische Details die man sich nicht merken muss, die aber trotzdem Pflicht sind zu sagen.
    \slidehead
    Es ist nicht ganz so einfach...

    \begin{itemize}[<+->]
        \item Der verwendete Typ muss immer zur Compile-Zeit fest sein
        \item Kann sehr oft weggelassen werden, weil der Compiler den Typ selber nachvollziehen kann\\
            \texttt{(Type inference)}
            %        \item Bei anderen Sprachen muss der verwendete Typ explizit festgelegt werden % Hier kann ich kein Beispiel geben, weil wir noch keine Variablen eingeführt haben
        \item[$\Rightarrow$] Das werdet ihr morgen sehen
    \end{itemize}

    % \pause % TODO: Ruben hilfe % Nein
    % In Kotlin gilt:
    % \begin{itemize}[<+->]
    %     \item \kotlininline{Double} kann nur Zahlen zwischen \kotlininline{-1.7976931348623157e+308} und \kotlininline{1.7976931348623157e+308} speichern
    %         \begin{itemize}
    %             \item Kleinster positiver \kotlininline{Double} ist \kotlininline{2.2250738585072014e-308}
    %         \end{itemize}
    %     \item \kotlininline{Int} kann nur Zahlen im Interval $[-2^{31}, 2^{31}-1]$ (inklusiv) speichern.
    %     \item Es gibt weitere Typen, die mehr speichern können
    % \end{itemize}
\end{frame}

%	\nextvid{Operatoren}{}

% !TeX document-id = {ecbf6175-be07-4601-b8eb-d3263d2bff07}
% !TeX TXS-program:compile = txs:///pdflatex/[--shell-escape]


%Deckblatt
% \subtitle{Kapitel 2: Spaß mit Daten}
% \titlegraphic{
% 	\begin{columns}
% 		\begin{column}{10cm}
% 			\vspace{4.32mm}
% 			\begin{center}
% 				{\huge Operatoren}
% 			\end{center}
% 			\begin{figure}
% 				\centering
% 			\includeinvertablegraphics[scale=.35]{media/too_old_for_this_shit.png}
% 				\\	\sffamily \tiny Bild: \href{https://xkcd.com/447/}{https://xkcd.com/447/}
% 			\end{figure}
% 		\end{column}
% 	\end{columns}}
% \maketitle

\section{Operatoren}
\begin{frame}
    \slidehead

    \begin{itemize}
        \item Symbol, das den \textbf{Compiler} zur Ausführung einer Aktion veranlasst
        \item Operatortypen:
    \end{itemize}
    \tikzstyle{rect} = [rectangle, rounded corners, minimum width=3cm, minimum height=1cm,text centered, draw, text width=5cm, fill=\IfDarkModeTF{red!30!black}{red!10}]
    \centering
    \begin{tikzpicture}[node distance=1.5cm]
        \node[rect] (1) {\textbf{Mathematische Operatoren}};
        \node[rect, below of=1] (2) {\textbf{Vergleichsoperatoren}};
        \node[rect, below of=2] (3) {\textbf{Bitoperatoren}};
    \end{tikzpicture}
\end{frame}

\subsection{Mathematische Operatoren}
\begin{frame}[fragile]
    \slidehead
    \vspace{-0.2cm}
    \tikzstyle{rect} = [rectangle, rounded corners, minimum width=3cm, minimum height=1cm, text centered, draw, text width=3cm, fill=\IfDarkModeTF{red!30!black}{red!10}]
    \vspace{0.05cm}
    \centering
    \begin{tikzpicture}
        \node[rect] (1) {Addition: \verb|+|};
        \node[rect, right=0.5em of 1] (2) {Subtraktion: \verb+-+};
        \node[rect, below=0.5em of 1] (3) {Multiplikation: \verb+*+};
        \node[rect, right=0.5em of 3] (4) {Division: \verb+/+};
        \node[rect] at ([yshift=-2em]$(3.south)!0.5!(4.south)$) (5) {Modulo: \verb+%+};
    \end{tikzpicture}
    \pause
    \vspace{-0.1cm}
    \begin{block}{Hinweis}
        \begin{itemize}
            \item Modulo liefert den \textit{Rest einer Division}: \kotlininline{23 % 4 = 3}
            \item Arithmetische Operationen zwischen Ganzzahlen und Fließkommazahlen resultieren in \textit{Fließkommazahlen}: \kotlininline{1 + 2.0 = 3.0} % Für die*den Programmieranfänger*innen mag es zwar total normal scheinen, dass man Kommazahlen mit Integers addieren kann, aber für den Computer ist das ein wichtiger Unterschied! Jetzt kommt der Moment wo ich alle überzeugen kann, dass der Computer wirklich komplett anders denkt als wir Menschen. Er hat ein komplett anderes Konzept, sich 0.1 vorzustellen.

                %Umfrage, glaubt ihr Computer können rechnen?
                \pause
            \item \kotlininline{0.1 + 0.2 =} \pause \kotlininline{0.30000000000000004} $\Rightarrow$ Fließkommazahlen sind nicht exakt
        \end{itemize}
    \end{block}
\end{frame}


\subsection{Vergleichsoperatoren}
\begin{frame}
    \slidehead

    \begin{itemize}
        \item Vergleichsoperatoren werten zu einem Wert vom Typ \kotlininline{Boolean} aus
        \item Beispiel:
            \kotlinfile{listings/vergleich.kts}
    \end{itemize}
\end{frame}


\subsection{Vergleichsoperatoren - Übersicht}
\begin{frame}
    \slidehead

    \begin{center}
        \vskip -7 pt
        \begingroup
        \fontsize{10pt}{11pt}\selectfont
        \begin{tabular}{|l|>{\ttfamily}c|>{\ttfamily}c|>{\ttfamily}c|}
            \hline
            \textbf{Operator}       & \textbf{Symbol} & \textbf{Beispiel}        & \textbf{Rückgabewert} \\ \hline
            gleich                  & ==              & \kotlininline{100 == 50} & \kotlininline{false}  \\
                                    &                 & \kotlininline{50 == 50}  & \kotlininline{true}   \\ \hline
            ungleich                & !=              & \kotlininline{100 != 50} & \kotlininline{true}   \\
                                    &                 & \kotlininline{50 != 50}  & \kotlininline{false}  \\ \hline
            größer als              & >               & \kotlininline{100 > 50}  & \kotlininline{true}   \\
                                    &                 & \kotlininline{50 > 50}   & \kotlininline{false}  \\ \hline
            größer als oder gleich  & >=              & \kotlininline{100 >= 50} & \kotlininline{true}   \\
                                    &                 & \kotlininline{50 >= 50}  & \kotlininline{true}   \\ \hline
            kleiner als             & <               & \kotlininline{100 < 50}  & \kotlininline{false}  \\
                                    &                 & \kotlininline{50 < 50}   & \kotlininline{false}  \\ \hline
            kleiner als oder gleich & <=              & \kotlininline{100 <= 50} & \kotlininline{false}  \\
                                    &                 & \kotlininline{50 <= 50}  & \kotlininline{true}   \\ \hline
        \end{tabular}
        \endgroup
    \end{center}

    %    \note{
    %        == geht auch für Strings!
    %    }
\end{frame}

\subsection{Bitoperatoren}
\begin{frame}
    \slidehead
    Ja, es gibt auch Bitoperatoren ...

    $\Rightarrow$ Die werden wir hier aber nicht behandeln
\end{frame}

\livecoding

\subsection{Operatorpräzedenz}
\begin{frame}
    \slidehead

    \kotlinfile{listings/rangfolge_basic.kts}
    \pause
    \kotlinfile{listings/rangfolge_parenthesis.kts}
    \pause
    \begin{itemize}
        \item Die Operatorpräzedenz gibt die \textbf{Reihenfolge} vor, in der Ausdrücke ausgewertet werden:
            \begin{enumerate}
                \item Arithmetische Ausdrücke
                    \begin{itemize}
                        \item[$\Rightarrow$] Wie aus Schulmathematik bekannt
                    \end{itemize}
                \item Vergleiche
                \item Logische Ausdrücke (hierzu später mehr)
                \item Zuweisungen
            \end{enumerate}
    \end{itemize}
\end{frame}

%\nextvid{Zeichenketten}{Operatoren Beispiel}
%\end{document}

% !TeX document-id = {ecbf6175-be07-4601-b8eb-d3263d2bff07}
% !TeX TXS-program:compile = txs:///pdflatex/[--shell-escape]

%
% %Deckblatt
% \subtitle{Kapitel 2: Spaß mit Daten}
% \titlegraphic{
% 	\begin{columns}
% 		\begin{column}{10cm}
% 			\begin{center}
% 				{\huge Zeichenketten}
% 			\end{center}
% 			\vspace{-1mm}
% 			\begin{figure}
% 				\centering
% 			\includeinvertablegraphics[scale=.35]{media/mu.png}
% 				\\	\sffamily \tiny Bild: \href{https://xkcd.com/815/}{https://xkcd.com/815}
% 				\\
% 				(Ich habe keinen passenden xkcd gefunden. Schlagt gerne welche vor.)
% 			\end{figure}
% 		\end{column}
% 	\end{columns}}
% \maketitle

\subsection{}
\section{Zeichenketten}

\begin{frame}
    \slidehead
    \begin{itemize}[<+->]
        \item Habt ihr vorhin kennengelernt
        \item Zeichenketten werden als \textit{Strings} bezeichnet
        \item Ein String kann beliebig viele Zeichen(\textit{Chars}) enthalten
        \item Strings können mit \textbf{+} konkateniert werden
    \end{itemize}
\end{frame}


%Vielleicht später mal nützlich:

%\begin{frame}
%	\slidehead
%	\begin{itemize}
%		\item String Slicing (Zerschneiden von Text)
%			\pythonfile{listings/String_Slicing_Listing.py}
%		\begin{block}{Wichtig}
%			\begin{itemize}
%				\item \textbf{[von : bis]} ist \textbf{inklusive} "`von"' und \textbf{exklusive} "`bis"'
%				\item Mathematisch: Ein Interval $[von, bis)$ bzw. $[von, bis[$
%				\item Der erste Index beginnt mit der Zahl \textbf{0}
%			\end{itemize}
%		\end{block}
%	\end{itemize}
%\end{frame}

%\begin{frame}
%	\slidehead
%	\begin{itemize}
%		\item String Slicing Teil 2
%			\pythonfile{listings/String_Slicing_Listing2.py}
%		\begin{block}{Wichtig}
%			\begin{itemize}
%				\item \textbf{[: bis]} ist \textbf{exklusive} und \textbf{[von :]} ist \textbf{inklusive} der Grenze
%				\item Auch hier beginnt der erste Index mit der Nummer \textbf{0}
%				\item Es können auch beide Grenzen weggelassen werden
%			\end{itemize}
%		\end{block}
%	\end{itemize}
%\end{frame}

\subsection{Funktionen und Operatoren}
\begin{frame}
    \slidehead

    \begin{itemize}
        \item Länge eines Strings
            \kotlinfile{listings/string_length.kts}

        \item Umwandeln in Kleinbuchstaben
            \kotlinfile{listings/string_lowercase.kts}

        \item Umwandeln in Großbuchstaben
            \kotlinfile{listings/string_uppercase.kts}
    \end{itemize}
\end{frame}


\subsection{Konvertierung}
\begin{frame}
    \slidehead
    \begin{block}{Warnung}
        Strings sind keine Zahlen
    \end{block}

    \kotlinfile{listings/string_not_a_number.kts}

    \pause

    \begin{itemize}[<+->]
        \item Strings müssen zum Rechnen konvertiert werden
        \item \textbf{.toInt()} konvertiert x zu \kotlininline{Int}
        \item \textbf{.toDouble()} konvertiert x zu \kotlininline{Double}
            %            \item \textbf{boolean(x)} konvertiert x zu boolean  % Das lassen wir weg, weil es durch leere Zeichenketten und Zahlen zu mehr Verwirrung sorgt als es hilfreich ist
        \item \textbf{.toString()} konvertiert x zu \kotlininline{String}
    \end{itemize}

\end{frame}

\section{Konvertierung}
\subsection{}
\begin{frame}
    \slidehead
    \begin{block}{Beispiel}
        \kotlinfile{listings/string_conversion.kts}
    \end{block}
\end{frame}

\livecoding

\section{Quiz}
\begin{frame}
    \slidehead
    \begin{itemize}[<+->]
        \item Wie kann ich \kotlininline{"} in Strings verwenden?
        \item Wo liegt der Unterschied zwischen \kotlininline{Int} und \kotlininline{Double}?
        \item In welcher Reihenfolge wird \kotlininline{2 + 4 * 2 == 9} ausgewertet?
        \item Was ist das Ergebnis?
        \item Wofür gibt es Kommentare?
    \end{itemize}
\end{frame}

%\nextvid{Variablen}{Zeichenketten Beispiel}

\end{document}
